\section{Extended Validation Details}
\label{app:extended-details}

\subsection{EVM Regression Details}
The Solidity fixture contains a state-effect target, a minimal relay, a
callback vault, and its callback client. All cases start from fresh local
Cancun EVM instances and are repeated three times. The following checks make
the representation and execution assumptions used by the evidence layer
observable.

\paragraph{Access lists.}
An EIP-2930 access list changes access accounting and transaction cost; it does
not prohibit reading an omitted storage slot~\cite{eip2930}. Two type-1
transactions to the same getter both succeed and return 42, with the empty
and populated-list receipts using \EmptyListGas{} and \PopulatedListGas{} gas.

\paragraph{Storage context and rollback.}
The dynamic-storage case derives a mapping slot from the key and mapping
position and verifies the raw value while slot zero remains unchanged. The
delegatecall case observes value 91 in caller storage and zero in the
implementation storage, connecting the effect to the delegatecall context
and ERC-1967 conventions~\cite{erc1967}. The transient-storage case observes
19 within one call and zero in a later call, while the rollback case writes
500 and reverts with the persistent value unchanged; the latter follows the
transaction-scoped semantics of EIP-1153~\cite{eip1153}.

\paragraph{Relay and callback behavior.}
An unauthorized relay caller fails without consuming the message, an
authorized caller succeeds, a blocked delivery remains pending until retry,
a duplicate fails, and an expired delivery leaves the destination unchanged.
In the vulnerable vault, a callback produces two logical payouts before the
credit is cleared; the corrected ordering produces one. These cases bind
caller, origin, retry key, deadline, receipt status, callback order, and
destination state to the observed effects.

\begin{table}[t]
\centering\small
\caption{Constructed Solidity regressions, each repeated three times. All
expected outcomes were observed.}
\label{tab:evm}
\begin{tabular}{@{}>{\raggedright\arraybackslash}p{0.36\columnwidth}>{\raggedright\arraybackslash}p{0.57\columnwidth}@{}}
\toprule
Case & Observed outcome\\
\midrule
Access list & Both reads succeed; gas differs\\
Dynamic slot & Derived mapping slot becomes 77\\
Delegatecall & Caller 91; implementation 0\\
Transient storage & Within-call 19; later call 0\\
Revert rollback & Failed receipt; stored value 0\\
Relay authorization & Wrong sender fails; governor succeeds\\
Relay retry & Failure preserves pending status; retry succeeds\\
Duplicate delivery & First succeeds; duplicate fails\\
Expiration & Late delivery fails without update\\
Decimal units & 0.75 represented as 750,000 at six decimals\\
Callback reentrancy & Logical payouts 2 versus 1 after correction\\
Prefix order & Equal final states; one unsafe prefix\\
\bottomrule
\end{tabular}
\end{table}

Slither analyzes the same four fixture contracts and reports 14 findings across
nine detector categories~\cite{feist2019slither,slitherrepo}. Its findings are
source-level warnings; BTG uses the same fixture to check ordered effects
against a joint property.

\begin{table}[t]
\centering\small
\caption{Slither 0.11.3 output on the fixture source.}
\label{tab:slither}
\begin{tabular}{@{}lr@{}}\toprule
Detector category & Findings\\\midrule
Controlled delegatecall & 1\\
Reentrancy without ETH transfer & 1\\
Benign reentrancy & 1\\
Unchecked low-level return & 1\\
Missing zero-address checks & 2\\
Timestamp use & 1\\
Assembly & 2\\
Low-level calls & 4\\
Immutable-state suggestion & 1\\
\bottomrule\end{tabular}
\end{table}

\subsection{Dependency-Graph Benchmark}
\label{app:dependency-overhead}
We measure materialized dependency construction separately from exploration.
The construction experiment has \DependencyConstructionFixtures{} fixtures:
32, 128, or 512 actions; zero, $n$, or $4n$ distinct pair predicates over
separate write slots; and a same-slot storage clique. Each configuration has
three fixed seeds. Seven warmed constructions per seed are timed; a separate
\texttt{tracemalloc} run records Python allocation peak. The measurement is
construction allocation, not whole-process resident memory.

With 512 actions and 2,048 pair predicates, construction takes
\DependencySparseMS{} ms with \DependencySparseMiB{} MiB peak allocation.
The 512-action storage clique contains 130,816 edges and takes
\DependencyDenseMS{} ms with \DependencyDenseMiB{} MiB peak. Thus edge
density, rather than action count alone, determines materialization cost.

The end-to-end experiment contains \DependencyExplorationFixtures{} fixtures
with 4, 6, or 8 actions, four predicate densities, and two seeds. It compares
on-demand dependency scans with cached adjacency using the same sleep-set DFS
and checks both relations against exhaustive enumeration. Across 328,512
exhaustive orders, the reduced runs retain 82,630 orders with zero relation or
joint-outcome mismatches. At zero property edges one order remains; when all
28 pairs are coupled all 40,320 orders remain. The raw per-seed runs and
generated tables are retained in the artifact.

\input{experiments/results/dependency_overhead_table.tex}

\subsection{Scale and Independent Cross-Check}
The scale study extends the deterministic overwrite generator from three--six
to three--eight labeled actions, producing 192 models and 1,479,360 complete
orders. Every method receives the identical model, initial state, and
predicate; only the observation boundary or dependency relation changes.
Property-aware POR retains 35,874 orders with zero disagreements against
exhaustive prefix checking. At eight actions this is a 98.43\% reduction;
the matched ablations report 106 disagreements for final-only checking, 192
for a fixed order, and 52 for storage-only POR.
\input{experiments/results/scale_ablation_table.tex}
\begin{figure*}[t]
\centering
\includegraphics[width=0.94\textwidth]{figures/verified_schedules.pdf}
\caption{Scale-study results generated from the saved schedule records. Left:
total complete orders for 32 constructed models at each action count; right:
unsafe models missed relative to exhaustive prefix checking.}
\label{fig:schedules}
\end{figure*}

The independent monitor exhausts zero to four labeled overwrite actions, three
binary-valued slots, all eight binary initial states, and the predicates
$x\geq y$ and $y\geq z$. This is
$8\sum_{n=0}^{4}6^n=12{,}440$ configurations and 259,832 exhaustive orders.
It compares joint pairs of terminal state and violated-predicate bitmask,
including the empty prefix. Property-aware reduction retains 177,592 orders
with zero outcome-set mismatches; storage-only reduction differs in 2,836
configurations. The complete aggregate and per-size output is
\path{experiments/results/crosscheck.json}.

\subsection{Extended Deployment and Coverage Protocol}
Protocol adapters preserve the immediate caller, logical origin, route,
bridge version, payload bytes, message identity, block hashes, deployed code
hashes, and relevant state. Unsupported selectors, upgrade patterns, missing
snapshots, and ambiguous intent are reported before a security decision.
Properties are tagged as official requirements, case-study policies, or
illustrative assumptions. Historical cases, injected faults, and synthetic
regressions remain separate evidence classes. The artifact guide below lists
the exact commands, manifests, source hashes, and raw outputs.
